\documentclass{article}
\usepackage{amsmath,amssymb,mathtools}
\usepackage[margin=1in]{geometry}
\allowdisplaybreaks
\usepackage[UTF8]{ctex}
\begin{document}

\section*{供应商信用合同优化}
\begin{equation}
\begin{aligned}
\max_{w_c,w_t}\quad \Pi_s
&=-cq^w+w_cq_c^w
 +(1-\alpha_t)\int_{\theta_b^w}^{\theta_t^w}(x-\theta_b^w)\,dF(x)
 +(\theta_t^w-\theta_b^w)\bar F(\theta_t^w),\\
\text{subject to}\quad &c\leq w_c\leq w_t\leq 1,\\
&(q^w,q_c^w,\theta_b^w,\theta_t^w)
 \text{ are the retailer's best response specified below.}
\end{aligned}\tag{7}
\end{equation}
最大化与呈现目标相结合，将第5节中的供应商目标与第3.1节中的合同领域结合起来，同时考虑其对命题1的依赖性。所选模型是完整的合同设计问题，包括没有外部融资情况下的响应；它并不局限于外部融资的表达式(8)。

\section*{固定的输入、决策和时间安排}
有一个供应商、一个零售商和一个产品。固定的单位生产成本为$c$。零售价、残值、无风险利率、零售商的初始现金以及违约成本的比例都符合要求。
\[
p=1,\qquad s=0,\qquad r_f=0,\qquad K>0,\qquad
\alpha_b,\alpha_t\in[0,1].
\]
需求分布是固定的，随机需求$\widetilde D$依赖于$[0,+\infty)$、实现方式$D$、累积分布函数$F$以及密度函数$f$。与之相关的函数有：
\[
\bar F(x)=1-F(x),\qquad h(x)=\frac{f(x)}{\bar F(x)},\qquad
 g(x)=xh(x).
\]
该企业的失败率呈上升趋势。未满足的需求被浪费了。零售商只能依靠未来的收入来偿还外部债务。除了违约成本之外，金融市场是完全竞争性的；各方都是风险中立的，拥有对称的信息，并且不存在道德风险。

供应商首先提供现金和信用价格$(w_c,w_t)$。然后，零售商选择总订单及其融资安排；交货和现金支付在需求实现之前进行，而还款则发生在需求实现之后。零售商的库存数量和分配情况如下：
\[
q\geq0,\qquad q_c\in[0,q],\qquad q_t=q-q_c.
\]
这里的$q_c$是指以现金价格购买的库存，不一定是通过内部现金完全融资获得的库存；$q_t$则是指通过商业信用购买的库存。现金支付为$w_cq_c$，而名义上的商业信用则为$w_tq_t$。提前付款的折扣为
\[
d_t=1-\frac{w_c}{w_t}.
\]
网络术语表示为$w_c=w_t$；两部分组成的术语表示为$w_c<w_t$。上标$w$表示零售商对固定合同的响应；上标$*$则表示内生最优合同及其所引发的各数量。

\paragraph{现金领域存在差异。}
第3节中提到了$K>0$，这里将其作为输入前提保留下来。附录A中列出了$K\geq0$，而命题1则列出了从零开始的区间。这些区间在下面被记录下来，同时也没有削弱第3节中的前提。

\section*{融资与零售商付款的定义}
银行贷款的本金和还款门槛是
\[
B=(w_cq_c-K)^+,\qquad \theta_b=(1+r_b)B,
\]
其中，$r_b$表示银行利率。银行贷款优先于商业信用。只有当$D\geq\theta_b$满足条件时，银行贷款才会被完全偿还；而当$D<\theta_t$发生时，商业信用就会违约。商业信用的门槛是
\begin{equation}
\theta_t=\theta_b+w_t(q-q_c).\tag{1}
\end{equation}
比例$\alpha_b$和$\alpha_t$是与违约时被没收的财产价值成正比的无用成本。如果比例为零，则相应的债权人不会遭受任何回收损失。

表2中列出的视界终点收益
\[
\begin{array}{c|cccc}
& D\in[0,\theta_b)&D\in[\theta_b,\theta_t)&D\in[\theta_t,q]&D>q\\ \hline
\text{Bank}&(1-\alpha_b)D&\theta_b&\theta_b&\theta_b\\
\text{Supplier}&0&(1-\alpha_t)(D-\theta_b)&\theta_t-\theta_b&\theta_t-\theta_b\\
\text{Retailer}&0&0&D-\theta_t&q-\theta_t
\end{array}
\]
这些是终端收益，而不是供应商的总利润，其中包括生产成本和初始现金销售。

\paragraph{现金并未耗尽：$w_cq_c<K$。}
零售商既不使用银行贷款，也不使用商业信用；$q_t=0$。零售商的目标是最大化……
\[
\Pi_r=\int_0^q\bar F(x)\,dx-w_cq.
\]

\paragraph{现金已耗尽：$w_cq_c\geq K$。}
具有竞争力的银行定价要求
\[
\int_0^{\theta_b}(1-\alpha_b)x\,dF(x)
 +\theta_b\bar F(\theta_b)=B=w_cq_c-K.
\]
定义
\[
\hat F_b(x):=\bar F(x)[1-\alpha_bg(x)].
\]
然后，银行的定价关系以及打印出来的利率表达式就是如此。
\begin{equation}
w_cq_c-K-\int_0^{\theta_b}\hat F_b(x)\,dx=0,\tag{2}
\end{equation}
\begin{equation}
r_b=\frac{\theta_b-B}{B}
=\frac{\theta_b}{\displaystyle\int_0^{\theta_b}\hat F_b(x)\,dx}-1.\tag{3}
\end{equation}
零售商选择$q_c$和$q$，相当于选择$q_c$和$q_t$，以最大化……
\begin{equation}
\Pi_r=\int_{\theta_t}^{q}\bar F(x)\,dx-K\tag{4}
\end{equation}
在符合(1)、(2)以及上述数量范围的前提下，在这个资金耗尽的情况下，方程(3)实际上是一个贷款定价表达式。该论文并未为$B=0$这一零本金情况下的数值提供任何数值。

\section*{完整的印刷版零售商回复：命题1}
如命题1开头所打印的那样定义，
\[
\kappa^{nb}(w_c)=w_c\bar F^{-1}(w_c),
\]
\[
\kappa_r^b(w_c,w_t)
=w_c\bar F^{-1}\!\left(
\frac{w_c}{1-\alpha_bg\bigl(\hat F_b^{-1}(w_c/w_t)\bigr)}
\right)
-\int_0^{\hat F_b^{-1}(w_c/w_t)}\hat F_b(x)\,dx,
\]
\[
C(\theta)=\frac{K+\displaystyle\int_0^{\theta}\hat F_b(x)\,dx}
{\hat F_b(\theta)}-\theta.
\]
这三个响应分支如下，其中$q_t^w=q^w-q_c^w$以及来自附录A的数量域保持不变。所有数学符号都按原顺序保留下来。

\paragraph{1. 无需外部融资。}
对于$K\in[\kappa^{nb}(w_c),+\infty)$，
\[
q^w=q_c^w=\bar F^{-1}(w_c),\qquad
\theta_b^w=\theta_t^w=0.
\]

\paragraph{2. 仅限银行贷款。}
对于$K\in[\kappa^b(w_c,w_t),\kappa^{nb}(w_c))$，
\[
q^w=q_c^w=\bar F^{-1}\!\left(
\frac{w_c}{1-\alpha_bg(\theta_b^w)}\right),
\qquad \theta_t^w=\theta_b^w,
\]
其中，阈值是由唯一确定的。
\begin{equation}
w_c\bar F^{-1}\!\left(\frac{w_c}{1-\alpha_bg(\theta_b^w)}\right)
-\int_0^{\theta_b^w}\hat F_b(x)\,dx=K.\tag{5}
\end{equation}

\paragraph{3. 银行贷款与贸易信贷。}
对于$K\in[0,\kappa^b(w_c,w_t))$，
\[
\theta_b^w=\hat F_b^{-1}(w_c/w_t),\qquad
q_c^w=\frac{K+\displaystyle\int_0^{\theta_b^w}\hat F_b(x)\,dx}{w_c},
\]
并且$(q^w,\theta_t^w)$共同满足
\[
w_t\bar F(\theta_t^w)=\bar F(q^w)
\]
以及
\begin{equation}
[\theta_t^w+C(\theta_b^w)]\bar F(\theta_t^w)
=q^w\bar F(q^w).\tag{6}
\end{equation}

\paragraph{打印出的阈值名称存在差异。}
在命题1中的定义中，下标位置为$\kappa_r^b$；而其分支区间及后续的证明则使用$\kappa^b$作为下标位置。不过，原文并未对后者进行单独的定义。两种形式都已被保留下来，而不是默默地将其视为相等。同样，第三分支的口头描述也没有排除零银行门槛的可能性：其印刷公式在净条件下显示为$\theta_b^w=\hat F_b^{-1}(1)$。

\paragraph{支持性响应条件。}
附录 C，引理 C.3 指出，对于给定的$(w_c,w_t)$，当且仅当满足 (2) 和 (1) 的条件下，包含$q_t^{w,3}>0$的响应$(q_c^{w,3},q_t^{w,3},\theta_b^{w,3},\theta_t^{w,3})$才能在$\Pi_r=\int_{\theta_t}^{q_c+q_t}\bar F(x)\,dx$的情况下全局最优。
\begin{equation}
\hat F_b(\theta_b^{w,3})=\frac{w_c}{w_t},\tag{69}
\end{equation}
\begin{equation}
w_t\bar F(\theta_t^{w,3})=\bar F(q_c^{w,3}+q_t^{w,3})\tag{70}
\end{equation}
这是第三个分支的辅助特征描述，并非另一个供应商优化问题。附录C指出，技术引理的证明可以根据请求提供；但这些证明在此处不可用，且尚未经过验证。

\section*{条件性阈值表示法}
当使用外部融资时，第5节中所述内容保持不变。
\[
w_cq_c^w=K+\int_0^{\theta_b^w}\hat F_b(x)\,dx
\]
要写出与之前相同的供应商利润。
\begin{equation}
\begin{aligned}
\Pi_s={}&K+\int_0^{\theta_t^w}\bar F(x)\,dx-cQ(\theta_b^w,\theta_t^w)\\
&-\left[\alpha_b\int_0^{\theta_b^w}x\,dF(x)
+\alpha_t\int_{\theta_b^w}^{\theta_t^w}(x-\theta_b^w)\,dF(x)\right].
\end{aligned}\tag{8}
\end{equation}
引文中的集合和数量定义来自推论 3。
\[
\Theta:=\left\{(\theta_b,\theta_t)\;\middle|\;
\begin{aligned}
&\theta_t\geq\theta_b\geq0,\\
&[\theta_t+C(\theta_b)]\bar F(\theta_t)
 \leq g^{-1}(1)\bar F\bigl(g^{-1}(1)\bigr)
\end{aligned}\right\},
\]
\[
(\theta_b^w,\theta_t^w)\in\Theta,
\qquad
q^w=Q(\theta_b^w,\theta_t^w)
:=\min\left\{q:q\bar F(q)
 =[\theta_t^w+C(\theta_b^w)]\bar F(\theta_t^w)\right\}.
\]
推论3声称存在一种一一对应的映射关系，使得价格满足$w_c\leq w_t\leq1$的条件。原始合同中的限制条件$c\leq w_c$仍然属于完整设计模型的一部分；它并未被包含在$\Theta$的定义中。证明选择较小的根值，因为(8)在$q$中会逐渐减小，而这一选择能够确保$g(q)\leq1$得到满足。

\paragraph{范围与映射上的差异。}
方程（8）仅在外部融资的情况下被打印出来，它并不替代（7）在无需外部融资情况下的回应。推论3中的无限制一对一声明不能包含命题1的第一个分支的全部内容，因为在那个分支中，不同的合同可以给出$(\theta_b^w,\theta_t^w)=(0,0)$。其证明也打印了$1-w_c/w_t=\hat F_b(\theta_b^w)$，这与命题1的第三个分支（69）以及下面的（11）相矛盾。这里的操作响应链保留了（69），而不是那个矛盾的证明表达式。对于$\Theta$或根约定则，也没有提供任何额外的限制条件来修正这些声明。

\section*{链接到最优性特征描述，而非额外的可行性约束条件}
方程式（9）至（13）是阈值分析的衍生结果，也是必要的最优性条件。它们并非完整设计问题的最终替代方案。这些方程式以原始顺序被记录下来，并不会被强制应用于每一个可行的合同方案。

\begin{equation}
\frac{\partial\Pi_s}{\partial\theta_t^w}
=\bar F(\theta_t^w)-c\frac{\partial Q}{\partial\theta_t^w}
-\alpha_t(\theta_t^w-\theta_b^w)f(\theta_t^w).\tag{9}
\end{equation}
\begin{equation}
\frac{\partial\Pi_s}{\partial\theta_b^w}
=-c\frac{\partial Q}{\partial\theta_b^w}
+\left\{[\alpha_t-\alpha_bg(\theta_b^w)]\bar F(\theta_b^w)
-\alpha_t\bar F(\theta_t^w)\right\}.\tag{10}
\end{equation}
引理 1 给出了最优价格的表达式
\begin{equation}
w_c^*=\frac{\hat F_b(\theta_b^*)\bar F(q^*)}{\bar F(\theta_t^*)},
\qquad
w_t^*=\frac{\bar F(q^*)}{\bar F(\theta_t^*)},\tag{11}
\end{equation}
其中，$q^*=Q(\theta_b^*,\theta_t^*)$以及打印出来的替代选项如下所示：
\begin{equation}
\theta_b^*=0,\qquad
\frac{c}{\bar F(q^*)[1-g(q^*)]}
=\frac{1-\alpha_tg(\theta_t^*)}{1-(\theta_t^*+K)h(\theta_t^*)}
\geq\frac{\alpha_t[1-\bar F(\theta_t^*)]}{f(0)K\bar F(\theta_t^*)},\tag{12}
\end{equation}
或者
\begin{equation}
\begin{gathered}
\theta_b^*>0,\\
\frac{c}{\bar F(q^*)[1-g(q^*)]}
=\frac{1-\alpha_t(\theta_t^*-\theta_b^*)h(\theta_t^*)}
{1-[\theta_t^*+C(\theta_b^*)]h(\theta_t^*)}
=\frac{[\alpha_t-\alpha_bg(\theta_b^*)]\bar F(\theta_b^*)
-\alpha_t\bar F(\theta_t^*)}{C'(\theta_b^*)\bar F(\theta_t^*)}.
\end{gathered}\tag{13}
\end{equation}
这里$C'$是之前定义的$C$的导数。该引理陈述的是必要性，而非一般充分性。其证明在物理PDF第37页中给出了严格条件$\partial\Pi_s/\partial\theta_b\vert_{(0,\theta_t^*)}<0$，并认为它与(12)中的弱不等式等价；(12)的内容被原样保留。在(12)中并未添加假设$f(0)>0$或零分母约定。

\paragraph{来源范围。}
下列各式依据所提供的 Yang 与 Birge 原文 \emph{商业信用、风险分担与库存融资组合} 的第 3 至 5 节、命题 1、推论 3、引理 1、表 2 及附录 A 至 C。 第5.2节的零售商参与扩展、仅价格比较问题以及经验回归并不属于本选编内容。

\end{document}
